\section{Two principles and the combinatorial deduction}
\label{sec:framework}

We first state the two analytic principles in the precise forms used to
prove Theorem~\ref{thm:main}. This also fixes the order of all parameters.
The proof of the Prediction Principle occupies
Sections~\ref{sec:arithmetic}--\ref{sec:prediction}, with the independent
missing-corner Lemma~\ref{lem:cube-corner} proved in
Section~\ref{sec:cube-limits}.
Sections~\ref{sec:rough-progressions} and~\ref{sec:cube-limits} develop the
nilsequence tools for the Alignment Principle, proved in
Section~\ref{sec:alignment}.

\subsection{Ordered blocks and admissible scales}

For nonempty subsets $A,B\subseteq[n]$,
$A<B$ means $\max A<\min B$; we write $A<i$ for $A<\{i\}$.
For a list of multiplicative quantities $u_1,\ldots,u_n$, put
$u_A=\prod_{j\in A}u_j$. In particular this convention applies to
$t_A,h_A,b_A$, and $x_A$. Empty products, when they occur, equal one.

A \emph{block} is a set $B=T\cup\{i\}$ with $\varnothing\ne T<i$.
We call $i$ its \emph{pivot} and $T$ its \emph{tail}. For such a block let
\begin{equation}\label{eq:adding-blocks}
 \mathcal E_B=
 \bigl\{P\cup\{j\}:\varnothing\ne P<T<j<i\bigr\}.
\end{equation}
These are the blocks whose values will be added at the anchor $B$.
A \emph{chain of length $m$} is a list
$B_d=T_d\cup\{i_d\}$, $d\in[m]$, satisfying
\begin{equation}\label{eq:block-chain}
 \varnothing\ne T_1<T_2<\cdots<T_m<i_1<\cdots<i_m.
\end{equation}
All blocks in a chain are disjoint. Moreover, if $k<d$, then
$B_k\in\mathcal E_{B_d}$. This last property is why we use this particular
ordering of tails and pivots.

All asymptotic parameters below are indexed by positive integers
$w\to\infty$, and
\[
 W=W(w)=\prod_{p\le w}p,
\]
where the product is over primes. An integer is \emph{$w$-smooth} if all
its prime factors are at most $w$. We say that a positive quantity $A$
\emph{dominates powers} of a quantity $B\ge2$ if
$A/B^C\to\infty$ for every fixed $C>0$. Every finite index set,
complexity bound and tolerance is held fixed in this convention unless
additional uniformity is stated. Sampling an interval always means
sampling its integer points.

Fix once and for all a nonprincipal ultrafilter $\mathcal U$ on the
positive integers $w$, and write $\lim_{\mathcal U}$ for its limit.
Every bounded real sequence has such a limit, and it agrees with the
ordinary limit when that exists. Every set belonging to $\mathcal U$
meets every sufficiently late range of $w$. These are the limit and
extraction properties used below.

\begin{definition}[Admissible parameters]\label{def:admissible}
For a fixed $n$, an admissible family consists of positive integers
$M,h_1,\ldots,h_n,H_1,\ldots,H_n$, positive powers of two
$X_1,\ldots,X_n$, and independent random variables $t_1,\ldots,t_n$,
all depending on $w$, with the following properties.
\begin{enumerate}
\item $M$ and each $h_i$ are $w$-smooth multiples of $W^w$.
For every block $B$, and also for every singleton $\{i\}$,
$h_B\le M$. For every $A\in\mathcal E_B$, the ratio $h_A/h_B$
is an integer multiple of $W^w$.
\item Each $H_i$ is a multiple of $M$ and dominates powers of
$2+M+\prod_{\ell<i}X_\ell$. The quantity $\log X_i$ dominates powers
of $H_i$.
\item The law of $t_i$ is the harmonic $W$-unit probability measure
\begin{equation}\label{eq:harmonic-law}
 \mu_i(y)=\frac{\1_{X_i\le y<X_i^2}\1_{(y,W)=1}}{Z_i y},
 \qquad
 Z_i=\sum_{\substack{X_i\le y<X_i^2\\(y,W)=1}}\frac1y.
\end{equation}
\end{enumerate}
\end{definition}

Fixed positive rational factors have only finitely many denominator
primes and fixed denominator valuations. Consequently they can be
absorbed by $W^w$ for all sufficiently large $w$. In particular, every
color argument used below is eventually a positive integer. For a chain
and fixed positive rational $a_1,\ldots,a_m$, the ratios
$h_{B_k}a_k/(h_{B_d}a_d)$, $k<d$, are eventually positive integers:
the pair belongs to \eqref{eq:adding-blocks}. These facts allow us to
use rational scale lists without ever evaluating the coloring on a
noninteger.

\subsection{Piecewise nilsequence models}

A \emph{nilsequence of step at most $s$} is a sequence
\begin{equation}\label{eq:nilsequence-definition}
 k\longmapsto F(g^k x),\qquad k\in\Z,
\end{equation}
where $G$ is a connected simply connected nilpotent Lie group of step
at most $s$, $\Gamma\subset G$ is a lattice, $x\in G/\Gamma$, $g\in G$,
and $F$ is a Lipschitz function on $G/\Gamma$.
Throughout this paper a family has \emph{bounded complexity} if its
nilmanifolds belong to a fixed finite list and its observables have
uniformly bounded supremum and Lipschitz norms for fixed smooth metrics.
No bound is imposed on the translating elements $g$ or the points $x$.
Constants are included. Products and Lipschitz combinations of any
fixed finite number of such families remain bounded-complexity
nilsequences of the same maximal step, by taking product nilmanifolds.

\begin{definition}[Piecewise models]\label{def:piecewise-model}
Fix finite sets $\mathcal A\subset\Q_{>0}$ and $[r]$, and an admissible
family. A system of step-at-most-$s$ models consists of functions
$S_{B,a,c}:\Z\to[0,1]$, one for each block $B$, $a\in\mathcal A$,
and $c\in[r]$. If $i=\max B$, then on each interval
$[kH_i,(k+1)H_i)$ and each residue class modulo $M$, the function is a
nilsequence in the progression index. The complexity is bounded
uniformly over $w$, all pieces, and all model indices. The representing
observables may be, and are, chosen $[0,1]$-valued on their entire
nilmanifolds.
\end{definition}

The sequence and its nilmanifold may change from piece to piece within
the fixed finite list. This freedom is necessary in the inverse-theorem
argument. Real parts followed by clipping to $[0,1]$ justify the final
sentence of Definition~\ref{def:piecewise-model} without increasing the
nilpotence step.

For a block $B=T\cup\{i\}$ define its divisor weight on all of $\Z$ by
\begin{equation}\label{eq:divisor-weight}
 \nu_B(y)=\E_{\sigma=t_T}\sigma\1_{\sigma\mid y}.
\end{equation}
The expectation uses the raw laws of the tail variables only.
For a fixed tail product $\sigma$, we have $(\sigma,W)=1$, so
the conditional mass of $\sigma t_i$ at $y$ is
\[
 \frac{\sigma\1_{\sigma\mid y}\1_{(y,W)=1}}{Z_i y}
 \quad\text{on }[\sigma X_i,\sigma X_i^2).
\]
Thus $\nu_B$ averages the divisibility factor inserted by multiplication
by the tail. Corollary~\ref{cor:product-law} controls the change of
endpoints back to $[X_i,X_i^2)$ and shows that the sum of the absolute
differences between the masses of $\operatorname{Law}(t_B)$ and
$\nu_B\mu_i$ tends to zero.
Whenever several copies of a weight are expanded, they receive
independent divisor samples. In particular, two occurrences of
$\nu_B$ do not share a sample unless this is expressly stated.

\subsection{Prediction and alignment}

We use finite lists $\mathcal B\subset\Q_{>0}^n$ and
$\mathcal A\subset\Q_{>0}$ satisfying
\begin{equation}\label{eq:scale-list-closure}
 b_B\in\mathcal A\quad\text{for all }b\in\mathcal B
 \text{ and all blocks }B.
\end{equation}
The coloring is a map $\chi:\N\to[r]$.

\begin{principle}[Prediction]\label{pr:prediction}
For every $m\ge2$ there is an integer $s=s(m)$ with the following
property. Given $n,r,\chi$, finite lists satisfying
\eqref{eq:scale-list-closure}, and tolerances
$0<\tau<1/4$, $\eta>0$, there are admissible parameters and
step-at-most-$s$ models such that the following conclusions hold along
the fixed ultrafilter $\mathcal U$ in $w$.
\begin{enumerate}
\item For each $B,a,c$,
\begin{equation}\label{eq:prediction-calibration}
 \lim_{\mathcal U}\PP\bigl(
 \chi(h_Bat_B)=c,\ S_{B,a,c}(t_B)\le2\tau\bigr)
 \le3\tau+\eta.
\end{equation}
\item Fix a chain $B_1,\ldots,B_m$, $b\in\mathcal B$ and a color
$c\in[r]$. Put
\[
 c_d=h_{B_d}b_{B_d},\qquad
 f_d(y)=\1_{\chi(c_dy)=c},\qquad
 \nu_d=\nu_{B_d},\qquad S_d=S_{B_d,b_{B_d},c}.
\]
Here $f_d$ is defined on positive integers and may be extended by zero
elsewhere. For $\varnothing\ne J\subseteq[m]$ write $d(J)=\max J$ and
\begin{align}
 L_J(z)&=\sum_{k\in J}\frac{c_k}{c_{d(J)}}z_k,
 \label{eq:sum-forms}\\
 U(z)&=\prod_{\varnothing\ne J\subseteq[m]}
 \1_{\chi(\prod_{k\in J}c_kz_k)=c}.
 \label{eq:product-mask}
\end{align}
Then the nonnegative weighted count
\begin{equation}\label{eq:weighted-count}
 \mathcal I_{b,c,\mathbf B}=
 \E_{z\sim\bigotimes_{d=1}^m\mu_{i_d}}
 U(z)\prod_{d=1}^m\nu_d(z_d)
 \prod_{\substack{J\subseteq[m]\\|J|\ge2}}
 (f_{d(J)}\nu_{d(J)})(L_J(z))
\end{equation}
satisfies, for $z(t)_d=t_{B_d}$,
\begin{equation}\label{eq:prediction-counting}
 \lim_{\mathcal U}
 \abs*{\mathcal I_{b,c,\mathbf B}
 -\E_t U(z(t))
 \prod_{\substack{J\subseteq[m]\\|J|\ge2}}
 S_{d(J)}(L_J(z(t)))}\le\eta.
\end{equation}
\end{enumerate}
\end{principle}

The step bound depends only on the requested chain length $m$.
The complexity bound may depend on all the fixed data and tolerances.
The extra weights on the sum forms in \eqref{eq:weighted-count} are
deliberate: they put the function at each sum under the same majorant
as the function at its last summand. Positivity of this weighted count
still gives an actual monochromatic configuration.

\begin{principle}[Alignment]\label{pr:alignment}
For every $n,r,s$ there are finite lists $\mathcal B,\mathcal A$
satisfying \eqref{eq:scale-list-closure} and a constant $\delta>0$,
depending only on $n,r,s$, such that the following holds.
For any admissible family, any system of step-at-most-$s$ models, and
any fixed $\tau>0$, define $\mathsf A_w$ to be the event that some
$b\in\mathcal B$ satisfies, simultaneously for every block $B$, color
$c\in[r]$, and subset $D\subseteq\mathcal E_B$,
\begin{equation}\label{eq:alignment-implication}
 \begin{split}
 S_{B,b_B,c}(t_B)>2\tau\quad\Longrightarrow\quad
 S_{B,b_B,c}\left(
 t_B+\sum_{A\in D}\frac{h_Ab_A}{h_Bb_B}t_A\right)>\tau.
 \end{split}
\end{equation}
Then
\begin{equation}\label{eq:alignment-probability}
 \liminf_{w\to\infty}\PP(\mathsf A_w)\ge\delta.
\end{equation}
In particular, $\delta$ is independent of $\tau$ and of the model
complexity bound.
\end{principle}

There is no assertion of a convergence rate uniform over all
complexities in \eqref{eq:alignment-probability}. Its positive lower
bound is uniform. This distinction allows us to choose the small
calibration tolerance before constructing the models.

\subsection{A finite combinatorial selection}

We use the finite sums theorem in its following finite form: for every
$m,r$ there is $F=F(m,r)$ such that every coloring of $[F]$ with $r$
colors contains positive integers $u_1,\ldots,u_m$ whose nonempty subset
sums all lie in $[F]$ and have one color. This is the finite sums theorem
of Folkman--Rado--Sanders \cite[Theorem~2 and Corollary~1.1]{Sanders};
it also follows from Hindman's finite sums theorem and compactness
\cite{Hindman,HindmanStrauss}. Indeed, if the finite form failed,
the finitely branching tree of bad finite colorings, ordered by restriction,
would have an infinite branch. Its limiting coloring of $\N$ would
contradict the finite sums theorem. We require no distinctness of
the auxiliary $u_d$ here.

The following selection uses the Ramsey-by-cardinality and finite-sums
argument of \cite[Section~5, Steps~II--IV]{Alweiss}; the ordered tails
and pivots are chosen to match the additive shifts in
\eqref{eq:adding-blocks}.

\begin{lemma}[Selection of a product chain]\label{lem:chain-selection}
For every $m,r$ there is $n=n(m,r)$ such that, for every coloring
$\chi:\N\to[r]$ and every list $x_1,\ldots,x_n\in\N$, a chain
$B_1,\ldots,B_m$ has all nonempty products of
$x_{B_1},\ldots,x_{B_m}$ in a single color class.
\end{lemma}

\begin{proof}
Choose $F$ as above and put $n_0=2F$. Repeated finite Ramsey gives
an $n$ such that every coloring of the nonempty index subsets of
$[n]$ has an $n_0$-element subset on which the color depends only on
cardinality, when the number of colors is $r$. More precisely, apply
the finite Ramsey theorem successively to the subset sizes
$1,\ldots,n_0$, choosing the successive ambient bounds backward.
Apply this to the coloring $A\mapsto\chi(x_A)$, and call the
resulting cardinality color $\kappa(q)$, $1\le q\le n_0$.

Color $a\in[F]$ by $\kappa(2a)$. Choose $u_1,\ldots,u_m$ by the
finite sums theorem and put $q_d=2u_d$. Then $q_d\ge2$,
$\sum_dq_d\le n_0$, and all nonempty subset sums of the $q_d$ have
the same $\kappa$-color. In the homogeneous ordered index set,
choose successive groups of $q_1-1,\ldots,q_m-1$ indices for
$T_1,\ldots,T_m$, followed by $m$ indices for the pivots
$i_1<\cdots<i_m$. This uses $\sum_dq_d$ indices and satisfies
\eqref{eq:block-chain}. A product of the corresponding block
products is $x_{\bigcup_{d\in J}B_d}$, whose index-set size is
$\sum_{d\in J}q_d$. All these sizes have the chosen color.
\end{proof}

\subsection{Deduction of the main theorem}

\begin{proof}[Proof of Theorem~\ref{thm:main}, assuming the two principles]
The case $m=1$ follows by choosing any integer $a_1>R$.
For $m\ge2$, fix $s=s(m)$ from
Principle~\ref{pr:prediction} and $n=n(m,r)$ from
Lemma~\ref{lem:chain-selection}. Apply Principle~\ref{pr:alignment}
to fix $\mathcal B,\mathcal A$ and $\delta>0$.
Let $C$ be the number of triples $(B,a,c)$ indexing calibration
events, and let $K$ be the number of triples consisting of a scale
vector, a chain of length $m$, and a color. These are fixed finite
numbers depending on data already chosen.

Choose $0<\tau<1/4$ and then $\eta>0$ so that
\begin{equation}\label{eq:outer-tolerances}
 3C\tau<\frac\delta4,\qquad
 C\eta<\frac\delta4,\qquad
 K\eta<\frac{\delta\tau^{2^m}}4.
\end{equation}
Now apply Principle~\ref{pr:prediction}. Define $\mathsf G_w$ to
be the event that none of its calibration failures occurs. The
union bound and \eqref{eq:prediction-calibration} give
\[
 \lim_{\mathcal U}\PP(\mathsf G_w^c)
 \le C(3\tau+\eta)<\frac\delta2.
\]
Since an ordinary lower limit bounds every ultrafilter limit,
\eqref{eq:alignment-probability} yields
\begin{equation}\label{eq:good-alignment-mass}
 \lim_{\mathcal U}\PP(\mathsf A_w\cap\mathsf G_w)
 \ge\frac\delta2.
\end{equation}

On this intersection choose a witnessing $b\in\mathcal B$ and
apply Lemma~\ref{lem:chain-selection} to
$x_i=h_ib_it_i$, which are positive integers for all sufficiently
large $w$. It supplies a chain and a color $c$ for which
$U(z(t))=1$. In particular, the singleton product masks say
$\chi(h_{B_d}b_{B_d}t_{B_d})=c$ for every $d$. Because
$\mathsf G_w$ holds, all centers $S_d(t_{B_d})$ exceed $2\tau$.

For a nonsingleton $J\subseteq[m]$, let $d=\max J$. Each
$B_k$, $k\in J\setminus\{d\}$, belongs to $\mathcal E_{B_d}$.
The corresponding instance of \eqref{eq:alignment-implication}
therefore gives
\[
 S_d\bigl(L_J(z(t))\bigr)>\tau.
\]
There are $q=2^m-m-1$ such factors. Thus at least one of the $K$
nonnegative model integrands in \eqref{eq:prediction-counting}
is greater than $\tau^q\ge\tau^{2^m}$ on
$\mathsf A_w\cap\mathsf G_w$. Summing expectations and using
\eqref{eq:good-alignment-mass} shows that their sum has ultrafilter
limit at least $\delta\tau^{2^m}/2$. It follows from
\eqref{eq:prediction-counting} and \eqref{eq:outer-tolerances} that
\begin{equation}\label{eq:positive-total-count}
 \lim_{\mathcal U}
 \sum_{b,c,\mathbf B}\mathcal I_{b,c,\mathbf B}>0.
\end{equation}
In particular, the finite sum is bounded below by a positive number
on a set belonging to $\mathcal U$. We did not need a fixed
choice of witnessing chain across different samples or values of $w$.

For some sufficiently large $w$, one weighted count is positive.
Since it is a finite nonnegative average, there is an actual tuple
$z$ in its support for which every factor is positive. Put
$a_d=c_dz_d$. The product mask gives every nonempty product of
the $a_d$ in color $c$, including their singleton values.
For every nonsingleton $J$, positivity of its color factor gives
\[
 \chi\left(c_{d(J)}L_J(z)\right)
 =\chi\left(\sum_{k\in J}a_k\right)=c.
\]
The extra divisor factors cannot create a false positive; they only
restrict which tuples can contribute to the count.

It remains to obtain the prescribed separation. The finite rational scale lists
give a constant $C_0$ such that $c_d\le C_0M$ for all these counts,
and $c_d\ge1$ once it is a positive integer. If $d<e$, then
\[
 a_d<C_0M X_{i_d}^2,\qquad a_e\ge X_{i_e}.
\]
For $e\ge2$, put $V_e=2+M+\prod_{j<i_e}X_j$.
There is a constant $C_1$, depending only on the fixed finite data,
such that every support tuple satisfies
\[
 \sum_{d<e}a_d+\prod_{d<e}a_d\le C_1 V_e^{3m}.
\]
Indeed each preceding $a_d$ is at most $C_0V_e^3$, and there are
at most $m-1$ such terms. Admissibility makes $X_{i_e}$ dominate
every fixed power of $V_e$. Therefore, for the prescribed $R,D$,
all support tuples at sufficiently large $w$ satisfy
\eqref{eq:separated-elements}; the first element exceeds $R$ since
$a_1\ge X_{i_1}\to\infty$. The positive-count set in
\eqref{eq:positive-total-count} belongs to the nonprincipal
ultrafilter and hence meets every sufficiently late range of $w$.
Choose $w$ in that intersection. This proves both monochromaticity
and separation, and completes Theorem~\ref{thm:main}.
\end{proof}

\begin{proof}[Proof of Corollary~\ref{cor:distinct-expressions}]
Apply Theorem~\ref{thm:main} with $R=2$ and $D=1$.
Every $a_j$ exceeds the sum and the product of all its predecessors,
and every element is at least two. If two subset sums were equal,
cancel their common terms. The greatest remaining index on either
side would exceed the sum on the other side, a contradiction.
For products, cancel common factors and compare the greatest
remaining factor with the product of all preceding factors.
An empty side after cancellation equals one, which is also smaller
than every remaining factor. Thus both subset maps are injective.

To compare a subset sum with a subset product, first compare their
greatest indices. If these differ, the expression with the larger
index exceeds the other expression. If both have greatest index $j$,
the sum lies in $[a_j,2a_j)$, whereas a nonsingleton product is at
least $2a_j$. A singleton product equals $a_j$, which equals a sum
with greatest index $j$ only for the singleton sum. Hence the
intersection is exactly $A$, giving the stated cardinalities.
\end{proof}

\begin{corollary}[Finite interval form with prescribed divisibility]
\label{cor:finite-interval}
Fix integers $r,m,q\ge1$ and real numbers $R\ge2$, $D\ge1$.
There exists $N$ such that every $r$-coloring of $[N]$ admits
$A=\{a_1<\cdots<a_m\}\subset q\N$ with
$\FS(A)\cup\FP(A)\subseteq[N]$ monochromatic and
\[
 a_1>R,\qquad
 a_d>R\left(\sum_{k<d}a_k+\prod_{k<d}a_k\right)^D
 \quad(2\le d\le m).
\]
The subset sums and products have the distinctness properties of
Corollary~\ref{cor:distinct-expressions}.
\end{corollary}

\begin{proof}
First color all of $\N$, refining the given color by the residue
modulo $q$. Apply Theorem~\ref{thm:main} to this finite coloring
with $\max(m,2)$ elements. Their common residue $v$ also equals
the residue of the sum of the first two elements, so $v=2v$
modulo $q$ and $v=0$. Retain the first $m$ elements. The
separation and monochromaticity persist, and the preceding proof
gives the distinctness assertions.

If no $N$ worked, the finite colorings avoiding such a configuration
would form a finitely branching tree, closed under restriction and
with a node at every depth. An infinite branch would color $\N$
without the configuration just proved to exist. Indeed each finite
configuration and all its sums and products lie in some initial
interval. This contradiction proves the finite interval assertion.
\end{proof}

This compactness argument gives a finite bound, but the proof does
not supply a numerical estimate for $N$.
